\documentclass[11pt,a4paper]{article}
\usepackage[margin=2.2cm]{geometry}
\usepackage{amsmath,amssymb}
\usepackage{booktabs,array,longtable}
\usepackage[hidelinks]{hyperref}
\usepackage{xcolor}
\usepackage{enumitem}
\setlist{nosep,leftmargin=1.4em}
\newcommand{\code}[1]{\texttt{\small #1}}
\newcommand{\where}[1]{\par\noindent\textbf{In the code:} #1\par}
\newcommand{\num}[1]{\par\noindent\textbf{Worked number (from the engine):} #1\par}
\newcommand{\ask}[1]{\par\noindent\textbf{Likely question:} \emph{#1}\par}
\newcommand{\notused}[1]{\par\noindent\textcolor{gray}{\textbf{Not used:} #1}\par}

\title{Course concepts in the Cp engine:\\ theory, derivations and where each is used}
\author{CL 202 Chemical Engineering Thermodynamics, Topic 5 (Specific heats of multiphase fluids)\\
TODO: group members}
\date{November 2026}

\begin{document}
\maketitle

\noindent
This report maps the material of the four lecture sets (\emph{Ch.~1 Definitions and basics},
\emph{Ch.~4 Property relationships for mixtures}, \emph{Equations of state for pure fluids},
\emph{Equations of state for mixtures}) to the Cp engine. For every concept it gives the result, the
derivation that can be asked on the board, the place in the code, and a number produced by the engine.
Concepts from the lectures that the engine does \emph{not} use are stated as such. All numbers were
computed with the engine (repository \code{cp-engine-multiphase}); slide numbers refer to the lecture PDFs.

\tableofcontents
\bigskip

\noindent\textbf{One idea runs through everything.} The engine never averages heat capacities directly.
It builds the enthalpy $H(T)$ of the whole system as a state function and differentiates it:
\begin{equation}
C_{p,\mathrm{app}} = \left(\frac{\partial H}{\partial T}\right)_{P,\,n_{\mathrm{total}}},\qquad
H=\sum_{\text{phases }p}\sum_i n_i^p H_i^p(T,P)+\sum_{\text{liquid }p} n^p H^{E,p}(T,\mathbf{x}^p)+H^R_{\text{gas}}.
\label{eq:main}
\end{equation}
Every lecture concept below enters through one of the terms of this equation or through the phase
amounts $n_i^p(T)$.

%=====================================================================================================
\section{Ch.~1: Definitions and basics}

\subsection{System, boundary, process type (slides 2--3)}
The user's description of the process is a statement about the boundary: \emph{adiabatic} (no heat
interaction), \emph{rigid} (constant volume), or movable at constant pressure. The engine turns these into
rules.
\where{\code{cp\_engine/rules.py} (rules \code{adiabatic\_no\_heat}, \code{first\_law\_constant\_P},
\code{first\_law\_constant\_V}); input key \code{process}.}

\subsection{Postulate I: $n+2$ properties fix the state (slide 7)}
For a closed simple system of $n$ components the stable equilibrium state is fixed by two independent
intensive properties plus the $n$ amounts. Route~A (the property route) therefore needs $T$, $P$ and the
composition; if one is missing, the solver's backward search lists which inputs would complete a route.
\where{\code{solver.py}, \code{Solver.missing\_suggestions}; Route~A rule in \code{rules.py}.}
\ask{Why does your engine refuse to give $C_p$ from a composition alone? (Postulate I: one more intensive
property, $T$, is needed; $P$ defaults to 1~atm and the engine says so.)}

\subsection{Postulate II, equilibrium states and quasi-static paths (slides 6, 8, 9)}
Derived properties ($U$, $H$, $S$) exist only for stable equilibrium states. The engine assumes that at
every temperature the system is in its stable equilibrium state, i.e.\ heating is quasi-static: the phase
amounts $n_i^p(T)$ are always the equilibrium amounts. Consequence: the apparent heat capacity is an
\emph{equilibrium} property. Supercooling (a PCM droplet that stays liquid below $T_m$) is a non-equilibrium
state and is not described.
\ask{Is your apparent Cp what a DSC measures on cooling? (On slow heating, approximately yes; on cooling, a
supercooled emulsion is not in equilibrium, so no.)}

\subsection{First law for a closed system (slides 19, 21) $\Rightarrow$ Route~B}
From $dE=\delta Q+\delta W$ with $E=U$ (no kinetic or potential energy):
\[
\Delta U = Q + W_s + W_{pv},\qquad W_{pv}=-\int P\,dV .
\]
At constant volume $W_{pv}=0$, so $\Delta U=Q+W_s$. At constant pressure $W_{pv}=-P\Delta V$ and, with
$H=U+PV$,
\[
\Delta H=\Delta U+P\Delta V = Q+W_s \qquad(\text{constant }P).
\]
Then the mean heat capacities are $\bar c_p=\Delta H/(m\Delta T)$ and $\bar c_V=\Delta U/(m\Delta T)$.
Sign convention as in the lectures: $Q>0$ is heat added to the system, $W>0$ is work done on it.
\where{\code{rules.py}: \code{\_dH\_constP}, \code{\_dU\_constV}, \code{cp\_from\_dH}.}
\num{$Q=83.7$~kJ heats 2~kg of water from 20 to 30~$^\circ$C at constant $P$, $W_s=0$:
$\bar c_p=83\,700/(2\times10)=4185$~J\,kg$^{-1}$K$^{-1}$; the engine returns 4185.0, and its own
property route gives 4184.8 (the two routes are compared automatically).}
\ask{If the stirrer does 5~kJ of work, what changes? ($\Delta H=Q+W_s$; the engine adds $W$.)}

\subsection{Heat capacity from entropy (slides 33--35) $\Rightarrow$ Route~C}
At constant $P$, $\delta Q_\mathrm{rev}=dH=C_p\,dT$ and $dS=\delta Q_\mathrm{rev}/T$, so
\[
C_p=T\left(\frac{\partial S}{\partial T}\right)_P ,\qquad
\Delta S=\int_{T_1}^{T_2}\frac{m\,c_p}{T}\,dT \;\overset{c_p\ \text{const}}{=}\; m\,c_p\ln\frac{T_2}{T_1}
\;\Rightarrow\; c_p=\frac{\Delta S}{m\ln(T_2/T_1)} .
\]
Because $S$ is a state function, $\Delta S$ between the two equilibrium states does not depend on whether
the actual heating was reversible; the formula needs only constant $P$ and an (approximately) constant
$c_p$.
\where{\code{rules.py}: \code{cp\_from\_dS}.}
\num{For 2~kg of water, 20$\to$30~$^\circ$C, $\Delta S=\int m c_p/T\,dT=280.74$~J/K; the entropy route
returns $c_p=4184.8$~J\,kg$^{-1}$K$^{-1}$.}

\subsection{$C_p-C_V$ (exact relation used by the $C_V$ route)}
With $dS=(C_V/T)dT+(\partial P/\partial T)_V dV$ (Maxwell relation $(\partial S/\partial V)_T=(\partial
P/\partial T)_V$), differentiating at constant $P$ and using the triple product rule
$(\partial P/\partial T)_V=\alpha/\kappa_T$:
\[
C_p-C_V=T\left(\frac{\partial P}{\partial T}\right)_V\left(\frac{\partial V}{\partial T}\right)_P
=\frac{T\,V\,\alpha^2}{\kappa_T},\qquad
\alpha=\frac1V\Big(\frac{\partial V}{\partial T}\Big)_P,\ \kappa_T=-\frac1V\Big(\frac{\partial V}{\partial P}\Big)_T .
\]
\where{\code{models/cv\_cp.py}; rule \code{cp\_from\_cv}.}
\num{Water at 25~$^\circ$C: $c_p=4184.5$, $c_V=c_p-Tv\alpha^2/\kappa_T=4141.1$~J\,kg$^{-1}$K$^{-1}$.
Validated against NIST $c_V$: 0.02\,\% AARD.}

\subsection{Thermodynamic temperature (slides 29--30)}
All temperatures are converted to kelvin on input; temperature \emph{differences} are not offset
(\code{units.py}).
\notused{heat engines, Carnot efficiency and the theorem of Clausius are not used directly; only their
consequence, the entropy as a state function, is. The open-system first law (slides 22--24) is not
implemented as a separate route; for a steady-flow calorimeter $\dot Q=\dot m\,\Delta h$ gives the same
$\bar c_p=\Delta h/\Delta T$ per unit mass.}

%=====================================================================================================
\section{Ch.~4: Property relationships for mixtures}

\subsection{Partial molar properties and Euler's theorem (slides 2--9)}
For an extensive property $\underline{B}(T,P,N_1,\dots,N_n)$, homogeneous of degree one in the $N_i$,
Euler's theorem gives
\[
\underline B=\sum_i N_i\bar B_i,\qquad \bar B_i=\left(\frac{\partial\underline B}{\partial N_i}\right)_{T,P,N_{j\neq i}},
\qquad B=\sum_i x_i\bar B_i .
\]
\emph{Proof (board):} $\underline B(\lambda N)=\lambda\underline B(N)$; differentiate with respect to
$\lambda$ and set $\lambda=1$: $\sum_i N_i\,\partial\underline B/\partial N_i=\underline B$.

From the molar property $B(x)$, eliminating $x_i$:
\[
\bar B_k=B+\left(\frac{\partial B}{\partial x_k}\right)_{T,P,x[k,i]}-\sum_{j\neq i}x_j\left(\frac{\partial
B}{\partial x_j}\right)_{T,P,x[j,i]} \qquad(\text{slide 8}).
\]
Applied to $C_p=T(\partial S/\partial T)_P$: $\bar C_{p,i}=T(\partial\bar S_i/\partial T)_{P,N}$
(slide~5), which the engine evaluates as $\bar C_{p,i}=C_{p,i}^{\text{pure}}+\bar C_{p,i}^{E}$.
\where{\code{models/partial\_molar.py}: \code{partial\_molar} (differentiates $N\,B(N/\sum N)$ with respect
to $N_i$) and \code{partial\_molar\_from\_x} (the slide-8 formula); both are coded and agree.
Every liquid-solution result reports \code{breakdown.extras["partial\_molar\_cp\_J\_per\_molK"]}.}
\num{Ethanol (0.3) + water (0.7) at 25~$^\circ$C: $\bar C_{p,\mathrm{water}}=90.22$,
$\bar C_{p,\mathrm{ethanol}}=138.85$~J\,mol$^{-1}$K$^{-1}$;
$0.3\times138.85+0.7\times90.22=104.81=C_p$ of the mixture (Euler).}

\subsection{Gibbs--Duhem (slides 10--11)}
Subtracting $d\underline B=\sum\bar B_i dN_i+\sum N_i d\bar B_i$ from the Gibbs equation for $d\underline B$:
\[
\sum_i N_i\,d\bar B_i=\Big(\frac{\partial\underline B}{\partial T}\Big)_{P,N}dT+\Big(\frac{\partial\underline
B}{\partial P}\Big)_{T,N}dP
\;\Rightarrow\;\sum_i x_i\Big(\frac{\partial\bar B_i}{\partial x_j}\Big)_{T,P}=0 .
\]
\where{\code{partial\_molar.gibbs\_duhem\_residual}; test in \code{tests/test\_eos\_vle.py}.}
\num{For the NRTL activity coefficients of ethanol + water, $\sum_i x_i\,\partial\ln\gamma_i/\partial x_1
=1.4\times10^{-11}$, i.e.\ zero: the model is thermodynamically consistent.}

\subsection{Mixing functions and the ideal solution (slides 13--18)}
With pure components at the same $T,P$ as reference states, the ideal solution has
\[
\Delta G^{ID}_{mix}=RT\sum x_i\ln x_i,\quad \Delta S^{ID}_{mix}=-R\sum x_i\ln x_i,\quad
\Delta H^{ID}_{mix}=0,\quad \Delta V^{ID}_{mix}=0 .
\]
Since $\Delta H_{mix}^{ID}=0$ at every $T$, its temperature derivative is zero too, so the ideal-solution
heat capacity is simply
\[
C_p^{ID}=\sum_i x_iC_{p,i}\quad(\text{molar basis}) \;\Longleftrightarrow\; c_p^{ID}=\sum_i w_ic_{p,i}\quad(\text{mass basis}).
\]
For a suspension the same argument over phases (no mixing between phases) gives mass weighting,
$c_p=\sum_p w_pc_{p,p}$; weighting by volume fraction (``Model~I'') is dimensionally inconsistent.
\where{\code{models/ideal.py}; \code{partial\_molar.mixing\_functions} (reported for every liquid solution);
\code{models/heterogeneous.py}.}
\num{Ethanol (0.3) + water (0.7), 25~$^\circ$C: $\Delta S^{ID}_{mix}=5.079$~J\,mol$^{-1}$K$^{-1}$,
$\Delta G_{mix}=-1360$~J/mol, $\Delta H_{mix}=-717$~J/mol, $\Delta S_{mix}=2.16$~J\,mol$^{-1}$K$^{-1}$
(smaller than ideal: the excess entropy is negative). Ideal $C_p=86.47$, real $C_p=104.81$~J\,mol$^{-1}$K$^{-1}$.}

\subsection{Excess functions, activity coefficients and Gibbs--Helmholtz (slides 24--33)}
$B^{EX}=B-B^{ID}$ at the same $T,P,x$, and $RT\ln\gamma_i=\bar G_i^{EX}=(\partial\,\underline{G}^{EX}/\partial
N_i)_{T,P,N_j}$. From $G=H-TS$ and $(\partial G/\partial T)_P=-S$ follows the Gibbs--Helmholtz equation,
\[
\left(\frac{\partial (G^{E}/RT)}{\partial T}\right)_{P,x}=-\frac{H^{E}}{RT^2}
\;\Rightarrow\;
H^E=-RT^2\frac{\partial g}{\partial T},\qquad
C_p^E=\left(\frac{\partial H^E}{\partial T}\right)_{P,x}=-R\left(2T\frac{\partial g}{\partial T}+T^2\frac{\partial^2
g}{\partial T^2}\right),
\]
with $g=G^E/RT$, and in partial molar form (slide~33) $(\partial\ln\gamma_i/\partial T)_{P,x}=-\bar H_i^{E}/RT^2$.
The mixture heat capacity is then $C_p=\sum x_iC_{p,i}+C_p^E$.

The lectures use the one-parameter Margules model $G^E/RT=C\,x_1x_2$ (slide 31) as the example. The engine
uses NRTL (Renon and Prausnitz, 1968), which is multicomponent and has temperature-dependent parameters,
\[
\frac{G^E}{RT}=\sum_i x_i\frac{\sum_j\tau_{ji}G_{ji}x_j}{\sum_k G_{ki}x_k},\qquad G_{ij}=\exp(-\alpha_{ij}\tau_{ij}),
\quad \tau_{ij}=a_{ij}+\frac{b_{ij}}T+e_{ij}\ln T .
\]
The temperature dependence matters: $C_p^E$ is a \emph{second} temperature derivative of $G^E$. Parameters
fitted to vapour--liquid equilibrium fix $G^E$ but not its temperature derivatives and give $C_p^E$
six or more times too small; the engine therefore uses parameters fitted to measured $H^E$.
\where{\code{models/nrtl.py}: \code{gE\_RT}, \code{ln\_gamma}, \code{hE}, \code{cpE} (central differences,
$h=10^{-4}T$); \code{partial\_molar.ln\_gamma\_from\_gE}, \code{partial\_molar\_hE}.}
\num{Ethanol (0.3) + water (0.7), 298.15~K: $H^E=-717$~J/mol, $C_p^E=+18.3$~J\,mol$^{-1}$K$^{-1}$.
$\ln\gamma$ from differentiating $N G^E$: $(-0.04701,\ 0.10913)$; analytic NRTL: $(-0.04701,\ 0.10913)$.
$\bar H^E_i$ from $-RT^2\partial\ln\gamma_i/\partial T$: $(-154,\ -958)$~J/mol, and
$0.3(-154)+0.7(-958)=-717=H^E$.}

\subsection{Fugacity, Lewis--Randall and Henry's law (slides 19--29)}
$\mu_i=RT\ln\hat f_i+\lambda_i(T)$; $\hat\phi_i=\hat f_i/(y_iP)$ for a gas; $\gamma_i=\hat f_i/(x_if_i)$ with
the pure liquid as reference (Lewis--Randall). These are used in the vapour--liquid equilibrium
(section~\ref{sec:vle}). The infinite-dilution activity coefficient $\gamma_i^\infty$ (the Henry's-law
reference, slides 28--29) is used to \emph{test} the NRTL parameters against independent measurements.
\where{\code{models/vle.py}; validation case~2 ($\gamma^\infty$ of ethanol and acetone in water).}

\subsection{Two class problems solved with the same relations}
\paragraph{Problem 9.13 (slide 35).} $RT\ln\gamma_B=(3800-8T)(1-x_B)^2$ is the two-suffix Margules form, so
$G^E=(3800-8T)\,x_Bx_C$ (J/mol). Then $S^E=-\partial G^E/\partial T=8x_Bx_C$ and
$H^E=G^E+TS^E=3800\,x_Bx_C$. Adding 1~mol benzene to 2~mol of solution with 80\,\% cyclohexane
(0.4~mol B, 1.6~mol C $\to$ 1.4~mol B, 1.6~mol C), isothermal at 300~K:
\[
\Delta H=3800\Big(\tfrac{1.4\times1.6}{3}-\tfrac{0.4\times1.6}{2}\Big)=1621.3\ \mathrm J,\qquad
\Delta S=\big[-R\textstyle\sum n_i\ln x_i+\underline S^E\big]_{\text{final}}-\big[\cdots\big]_{\text{initial}}=12.33\ \mathrm{J/K}.
\]
\paragraph{Problem 9.24 (slide 36).} With $\Delta\underline H=[100N_1N_2+50(N_1N_3+N_2N_3)]/(N_1+2N_2+0.5N_3)$
(homogeneous of degree one), at $x_1=0.2$, $x_2=0.5$:
\[
\overline{\Delta H}_2=\frac{(100x_1+50x_3)D-2\,N}{D^2}=\frac{35\times1.35-2\times20.5}{1.35^2}=3.429\ \mathrm{J/mol},
\]
with $N=20.5$, $D=1.35$. The engine's \code{partial\_molar} and the slide-8 formula both give 3.4294.
For part (b), mixing equal amounts of $(\tfrac13,\tfrac23,0)$ and $(\tfrac13,0,\tfrac23)$ to the equimolar
ternary at steady state: $Q=\Delta H=19.048-\tfrac12(13.333+16.667)=+4.05$~J per mol of product (with
exact thirds), i.e.\ heat must be \emph{supplied}: the ``cooling load'' is $-4.05$~J/mol.

%=====================================================================================================
\section{Equations of state for pure fluids}

\subsection{Why an equation of state is needed for $C_p$}
The database gives ideal-gas heat capacities $C_p^{ig}(T)$. A real gas at pressure differs by the
departure (residual) function. From $dH=T\,dS+V\,dP$ and the Maxwell relation $(\partial S/\partial P)_T=-(\partial
V/\partial T)_P$:
\[
\left(\frac{\partial H}{\partial P}\right)_T=V-T\left(\frac{\partial V}{\partial T}\right)_P,\qquad
\left(\frac{\partial C_p}{\partial P}\right)_T=-T\left(\frac{\partial^2V}{\partial T^2}\right)_P .
\]
Integrating from $P=0$ (ideal gas) at constant $T$ gives the departure functions
$H^R=\int_0^P[V-T(\partial V/\partial T)_P]\,dP$ and $C_p=C_p^{ig}+C_p^R$, $C_p^R=(\partial H^R/\partial T)_P$.
For an ideal gas $V=RT/P$ and every residual is zero. For liquids the engine applies the same relation with the
Tait equation for $V(T,P)$ (a liquid correlation, not in the slides): water at 500~bar gives 4056
(NIST 4060)~J\,kg$^{-1}$K$^{-1}$.
\where{\code{models/eos.py} (gases); \code{database.py}: \code{cp\_pressure\_correction} (liquids, integral by
8-point Gauss--Legendre quadrature).}

\subsection{Virial equation (slides 5--8)}
Truncated pressure series $Z=1+BP/RT$, i.e.\ $V=RT/P+B(T)$. Then $V-T(\partial V/\partial T)_P=B-TB'$ does not
depend on $P$, so
\[
H^R=P\,(B-TB'),\qquad C_p^R=\Big(\frac{\partial H^R}{\partial T}\Big)_P=P(B'-B'-TB'')=-T\,P\,\frac{d^2B}{dT^2},
\qquad \ln\phi=\frac{BP}{RT}.
\]
$B$ from the Pitzer/Abbott correlation, $BP_c/RT_c=B^0+\omega B^1$, $B^0=0.083-0.422/T_r^{1.6}$,
$B^1=0.139-0.172/T_r^{4.2}$; the derivatives $dB^0/dT_r=0.6752\,T_r^{-2.6}$ etc.\ are coded analytically.
\where{\code{eos.py}: \code{VirialEOS.\_B}, \code{H\_residual}, \code{Cp\_residual}.}
\num{CO$_2$ at 350~K (slide 30 constants): $B=-83.73$~cm$^3$/mol (slide: $-83.73$).}

\subsection{Generic cubic equation (slides 11--17)}
\[
P=\frac{RT}{V-b}-\frac{a(T)}{(V+\varepsilon b)(V+\sigma b)},\quad a=\Omega_a\frac{R^2T_c^2}{P_c}\alpha(T_r,\omega),\quad
b=\Omega_b\frac{RT_c}{P_c},
\]
with the $(\varepsilon,\sigma,\Omega_a,\Omega_b)$ of van der Waals, RK, SRK and PR from slide~16. The cubic in
$Z$ is solved for its real roots; the largest is the vapour, the smallest the liquid, the middle one is
discarded ($(\partial P/\partial V)_T>0$, slide~17). With $\beta=bP/RT$, $q=a/(bRT)$ and
$I=\ln[(Z+\sigma\beta)/(Z+\varepsilon\beta)]/(\sigma-\varepsilon)$:
\[
\ln\phi=Z-1-\ln(Z-\beta)-qI,\qquad \frac{H^R}{RT}=Z-1+\Big(\frac{d\ln a}{d\ln T}-1\Big)qI .
\]
The vapour pressure follows from equal fugacities, $\phi^L=\phi^V$ (iterated $P\leftarrow
P\,\phi^L/\phi^V$).
\where{\code{eos.py}: \code{CubicEOS.Z\_roots}, \code{ln\_phi}, \code{H\_residual}, \code{psat}.}
\num{CO$_2$ at 280~K, $P^\mathrm{sat}$ (bar) / $V^L$ (cm$^3$/mol): vdW 52.55 / 81.14, RK 45.61 / 61.03,
SRK 41.98 / 58.42, PR 41.60 / 51.68; slide 33: 52.55 / 81.13, 45.62 / 61.03, 41.99 / 58.42, 41.60 / 51.68
(agreement within rounding). On our own five liquids (NIST data, $0.55\le T_r\le0.95$) the vapour-pressure
error is 224\,\% (vdW), 58\,\% (RK), 2.3\,\% (SRK), 1.4\,\% (PR): Soave's idea of fitting $\alpha(T)$ to
vapour pressures (slide 14), seen on real data.}

\subsection{Acentric factor and corresponding states (slides 2, 8)}
$\omega=-1-\log_{10}[P^\mathrm{sat}(T_r=0.7)/P_c]$. The engine computes $\omega$ from this definition with NIST
saturation data and takes $T_c$, $P_c$ from the NIST reference equations of state:
water 0.3443, methanol 0.5647, ethanol 0.6429, $n$-hexane 0.3003, toluene 0.2657.
\where{\code{tools/fit\_vapour\_pressure.py}; \code{data/components.json} (\code{critical}).}

\subsection{Real-gas heat capacity: worked example}
\num{Superheated steam, 500~$^\circ$C, 100~bar: ideal gas 2131, virial 2469, Peng--Robinson 2496,
NIST 2583~J\,kg$^{-1}$K$^{-1}$. PR departure: $H^R=-2317$~J/mol, $C_p^R=+6.58$~J\,mol$^{-1}$K$^{-1}$.
Over 189 NIST points (4 vapours, 0.5--10~MPa) the error falls from 12.1\,\% (ideal gas) to 5.3\,\% (PR).}
\notused{the multiparameter and SAFT equations (slides 21--29) are not implemented; the NIST reference data
used for validation are themselves generated by multiparameter reference equations.}

%=====================================================================================================
\section{Equations of state for mixtures}

\subsection{Mixing and combining rules (slides 4, 7, 9, 11)}
Virial: $B_{mix}=\sum_i\sum_jy_iy_jB_{ij}$ (exact), with the cross coefficient from the Prausnitz combining
rules $T_{c,ij}=\sqrt{T_{c,i}T_{c,j}}(1-k_{ij})$, $V_{c,ij}=[(V_{c,i}^{1/3}+V_{c,j}^{1/3})/2]^3$,
$Z_{c,ij}=(Z_{c,i}+Z_{c,j})/2$, $\omega_{ij}=(\omega_i+\omega_j)/2$, $P_{c,ij}=Z_{c,ij}RT_{c,ij}/V_{c,ij}$.
Cubic: $a=\sum\sum x_ix_j\sqrt{a_ia_j}(1-k_{ij})$, $b=\sum x_ib_i$ (van der Waals one-fluid rules).
\num{Equimolar CO$_2$ + CH$_4$ at 350~K: $B_{12}=-46.98$~cm$^3$/mol (slide 23: $-47.0$); $Z$ at 20~bar:
virial 0.9650, PR ($k_{12}=0$) 0.9569, PR ($k_{12}=0.10$) 0.9600 (slide 24: 0.9650, 0.9569, 0.9600).}

\subsection{Fugacity coefficients in a mixture (slides 5, 8, 12)}
Virial: $\ln\hat\phi_i=\frac{P}{RT}\big(2\sum_jy_jB_{ij}-B_{mix}\big)$; $\delta_{12}=0$ gives the
Lewis--Randall rule. Cubic: $\ln\hat\phi_i=\frac{b_i}{b}(Z-1)-\ln(Z-\beta)-\bar q_iI$,
$\bar q_i=q\,(2\sum_jx_ja_{ij}/a-b_i/b)$; for PR ($\sigma-\varepsilon=2\sqrt2$, $q=A/B$) this is exactly the
expression on slide~12.
\num{$\hat\phi_{\mathrm{CO_2}}$ at 350~K, 20~bar, equimolar (virial): 0.947 (slide 24: 0.947).}

\subsection{Phase equilibrium from one equation: bubble pressure (slides 13, 25--27)}
$x_i\hat\phi_i^L=y_i\hat\phi_i^V$, $K_i=\hat\phi_i^L/\hat\phi_i^V$; start from Wilson $K$-values, update
$y_i=K_ix_i$, scale $P$ by $\sum K_ix_i$ until $\sum y_i=1$.
\where{\code{eos.py}: \code{CubicEOS.bubble\_P}.}
\num{CO$_2$ + CH$_4$, 273.15~K, 6.375~MPa: liquid CH$_4$ fraction 0.186 with $k_{12}=0$ (slide 26: 18.6\,\%)
and 0.128 with $k_{12}=0.10$ (measured: 12.5\,\%). The parameter that moved $Z$ by 0.3\,\% moves the liquid
composition by about 50\,\%.}

\subsection{Liquid mixtures that boil: the $\gamma$--$\phi$ approach}\label{sec:vle}
For water + alcohols no constant $k_{ij}$ fits both phases (slide~15), so the engine uses a liquid-solution
model for the liquid and the equation of state only for the vapour. Equal fugacities with the
Lewis--Randall reference for the liquid:
\[
y_i\hat\phi_iP=x_i\gamma_i\,\phi_i^{sat}P_i^{sat}\qquad\Rightarrow\qquad
K_i=\frac{y_i}{x_i}=\frac{\gamma_iP_i^{sat}\phi_i^{sat}}{\hat\phi_iP}\quad(\text{Poynting factor}\approx1),
\]
with $\gamma_i$ from NRTL, $P_i^{sat}$ from Antoine equations and $\hat\phi_i$, $\phi_i^{sat}$ from the same
Peng--Robinson equation that gives the vapour's $H^R$. Bubble point: $\sum K_ix_i=1$; dew point:
$\sum y_i/K_i=1$. Between them, the isothermal flash: from the material balance $z_i=(1-V)x_i+Vy_i$ and
$y_i=K_ix_i$,
\[
x_i=\frac{z_i}{1+V(K_i-1)},\qquad
f(V)=\sum_i\frac{z_i(K_i-1)}{1+V(K_i-1)}=0 \quad(\text{Rachford--Rice}),
\]
and $f'(V)=-\sum z_i(K_i-1)^2/[1+V(K_i-1)]^2<0$, so the root in $0<V<1$ is unique.
\where{\code{models/vle.py}: \code{bubble\_T}, \code{dew\_T}, \code{flash}; \code{system.py} decides from the
bubble and dew points whether a mixture is liquid, boiling or vapour.}
\num{Water + ethanol at 1~atm: 5~mol\,\% ethanol boils at 363.5~K (90.3~$^\circ$C), so it is liquid at
85~$^\circ$C although pure ethanol boils at 78~$^\circ$C; the azeotrope is found at $x=0.90$ where bubble and dew
points meet (351.28~K). Flash of $z=0.3$ at 355~K: $V=0.114$, $x_{\mathrm{EtOH}}=0.265$,
$y_{\mathrm{EtOH}}=0.570$. Bubble points of 409 measured points from 24 papers: 0.26~K (ethanol) and
0.47~K (methanol) mean error.}
\notused{Kay's rule (slide 3), Huron--Vidal and Wong--Sandler mixing rules (slide 15), SAFT and multi-fluid
Helmholtz models for mixtures (slides 18--22).}

%=====================================================================================================
\section{The calculus and numerics behind the engine}
\begin{itemize}
\item \textbf{Chain rule for the apparent heat capacity.} With $H=H(T,\mathbf n(T))$:
$\dfrac{dH}{dT}=\underbrace{\sum n_i^pC_{p,i}^p}_{\text{ideal}}+\underbrace{\sum n^pC_p^{E,p}}_{\text{excess}}
+\underbrace{\sum_{i,p}\frac{\partial H}{\partial n_i^p}\frac{dn_i^p}{dT}}_{\text{latent}}$.
Latent heat enters only through the last term, so it is counted once
(\code{system.py}: \code{cp\_frozen}, \code{cp\_point}).
\item \textbf{Kirchhoff's law (path independence).} $H_L(T)=\int_{T_m}^TC_{p,L}\,dT$,
$H_S(T)=-\Delta H_{fus}+\int_{T_m}^TC_{p,S}\,dT$, so $H_L-H_S=\Delta H_{fus}+\int_{T_m}^T\Delta C_p\,dT$: the
latent heat released below $T_m$ is smaller when $\Delta C_p>0$ (\code{models/pure.py}).
\item \textbf{Solid--liquid equilibrium.} $\ln(x_k\gamma_k)=-\frac{\Delta H_{fus}}{R}\big(\frac1T-\frac1{T_m}\big)
+\frac{\Delta C_p}{R}\big(\frac{T_m}T-1-\ln\frac{T_m}T\big)$ from $\mu_k^{solid}=\mu_k^{liquid}$ and
Gibbs--Helmholtz with a constant $\Delta C_p$ (\code{models/phase\_change.py}).
\item \textbf{Smoothed melting of a pure component.} Melted fraction
$f(T)=\tfrac12[1+\mathrm{erf}((T-T_m)/\sigma\sqrt2)]$, the cumulative distribution of a Gaussian; $f'(T)$ has
area exactly 1, so the area of the latent peak is $n\,\Delta H_{fus}$ (plus a $\tfrac12\sigma^2\,d\Delta
C_p/dT$ correction). The width is a modelling choice; interval means are independent of it.
\item \textbf{Central differences.} $f'(x)\approx[f(x+h)-f(x-h)]/2h$ with error $h^2f'''/6$ (Taylor series);
used for $C_p^E$, the latent term and $C_p^R$. Checked: the finite difference of $H(T)$ converges to the
reported $C_p$ as $\Delta T^2$.
\item \textbf{Gauss--Legendre quadrature with a change of variable.} $\int_{P_0}^{P}f\,dp=\frac{P-P_0}{2}\sum_k
w_kf\big(\tfrac{P+P_0}2+\tfrac{P-P_0}2\xi_k\big)$: the factor $(P-P_0)/2$ is the Jacobian. Eight nodes reproduce
the closed-form Tait integral to $10^{-12}$ (relative).
\item \textbf{Energy conservation as a check.} $\int_{T_1}^{T_2}C_{p,\mathrm{app}}\,dT=H(T_2)-H(T_1)$ holds to
$10^{-7}$ or better through melting, crystallisation and boiling.
\end{itemize}

%=====================================================================================================
\section{Summary: concept $\to$ code $\to$ check}
\begin{longtable}{>{\raggedright}p{4.6cm} >{\raggedright}p{5.0cm} >{\raggedright\arraybackslash}p{5.6cm}}
\toprule
Lecture concept & Where in the code & How it is checked\\
\midrule\endhead
First law, $\Delta H=Q+W_s$, $\Delta U=Q+W$ & \code{rules.py} & Route~B equals Route~A to $<10^{-6}$\,\% (7 systems)\\
$C_p=T(\partial S/\partial T)_P$ & \code{rules.py} (\code{cp\_from\_dS}) & Entropy route gives 4184.8 for water\\
$C_p-C_V=Tv\alpha^2/\kappa_T$ & \code{models/cv\_cp.py} & NIST $c_V$: 0.02\,\%\\
Postulate I ($n+2$) & \code{solver.py} (missing inputs) & Tests in \code{tests/test\_solver.py}\\
Partial molar properties, Euler & \code{models/partial\_molar.py} & $\sum x_i\bar C_{p,i}=C_p$\\
Gibbs--Duhem & \code{partial\_molar.py} & Residual $10^{-11}$\\
Ideal solution, mixing functions & \code{models/ideal.py}, \code{partial\_molar.py} & Exact reduction tests\\
$G^E$, $\gamma$, Gibbs--Helmholtz, $C_p^E$ & \code{models/nrtl.py} & Measured $H^E$ and $C_p$ (held-out tests)\\
Fugacity, Lewis--Randall, Henry & \code{models/vle.py} & Bubble points; $\gamma^\infty$\\
Virial EOS, Abbott $B$ & \code{models/eos.py} & Slide 30 value reproduced\\
Cubic EOS, roots, $P^{sat}$ & \code{models/eos.py} & Slide 33 table reproduced\\
Departure functions $H^R$, $C_p^R$ & \code{models/eos.py}, \code{system.py} & NIST vapour $c_p$; two-route check of $H^R$\\
Mixing/combining rules, $\hat\phi_i$ & \code{models/eos.py} & Slides 23--24 reproduced\\
$\phi$--$\phi$ bubble pressure, $k_{ij}$ & \code{CubicEOS.bubble\_P} & Slide 26 reproduced\\
$\gamma$--$\phi$ VLE, Rachford--Rice & \code{models/vle.py} & 409 measured bubble points\\
\bottomrule
\end{longtable}

\section{Questions to prepare}
\begin{enumerate}
\item Derive $\Delta H=Q+W_s$ at constant pressure from the first law. Which work term disappears, and why?
\item Why is the apparent heat capacity defined as $dH/dT$ and not as a weighted average of $C_p$'s?
\item Show that an ideal solution has $C_p=\sum x_iC_{p,i}$.
\item Prove Euler's theorem for partial molar properties; then show that $\sum x_i\bar C_{p,i}=C_p$.
\item Derive $H^E=-RT^2\,\partial(G^E/RT)/\partial T$ and $C_p^E$ from it. Why do VLE-fitted parameters give a
poor $C_p^E$?
\item Derive $(\partial C_p/\partial P)_T=-T(\partial^2V/\partial T^2)_P$. What is it for an ideal gas?
\item With $Z=1+BP/RT$, derive $C_p^R=-TP\,d^2B/dT^2$.
\item Which root of the cubic is the vapour? Why is the middle root discarded?
\item Why do vdW and RK fail for vapour pressures while SRK and PR do not?
\item What is the Lewis--Randall rule, and when does $\delta_{12}=0$ reproduce it?
\item Why $\gamma$--$\phi$ and not $\phi$--$\phi$ for water + ethanol?
\item Derive the Rachford--Rice equation. Why is its root unique?
\item A 5\,\% ethanol solution at 85~$^\circ$C: liquid or vapour? How do you know?
\item Which postulate tells you how many inputs Route~A needs?
\item Is your apparent heat capacity an equilibrium property? What does that exclude?
\end{enumerate}

\end{document}
